% Part: incompleteness
% Chapter: representability-in-q
% Section: beta-function

\documentclass[../../../include/open-logic-section]{subfiles}

\begin{document}

\olfileid{inc}{req}{bet}
\olsection{لم تابع بتا}


برای نشان‌دادن اینکه در صورت دراختیارداشتن جمع، ضرب و~$\Char{=}$
می‌توانیم بازگشت اولیه را انجام دهیم، باید توابعی برای کار با دنباله‌ها
پرورش دهیم. (اگر توان‌رسانی را نیز داشتیم، کارمان آسان‌تر بود.) هنگامی
که بازگشت اولیه در اختیارمان بود، می‌توانستیم چیزهایی مانند «عدد اول
$n$اُم» را تعریف و رمزگذاری نسبتاً سرراستی انتخاب کنیم. اما اینجا
بازگشت اولیه نداریم---در واقع می‌خواهیم نشان دهیم بازگشت اولیه را
می‌توان با کمینه‌سازی انجام داد---پس باید هوشمندانه‌تر عمل کنیم.

\begin{lem}
\ollabel{lem:beta}
تابعی چون~$\beta(d,i)$ وجود دارد به‌گونه‌ای که برای هر دنبالهٔ
$a_0$،~\dots،~$a_n$ عددی چون~$d$ وجود داشته باشد که برای هر
$i\le n$، $\beta(d,i)=a_i$. افزون بر این، می‌توان~$\beta$ را تنها با
ترکیب و کمینه‌سازی منظم از توابع پایه تعریف کرد.
\end{lem}

$d$ را رمزی برای دنبالهٔ~$\tuple{a_0,\dots,a_n}$ در نظر بگیرید که
$\beta(d,i)$ عضو~$i$اُم را بازمی‌گرداند. (توجه کنید این «رمزگذاری»
از رمزگذاریِ توان‌های اعداد اول که پیش‌تر با آن آشنا شدیم
\emph{استفاده نمی‌کند}!) ادعای لم بسیار کمینه است: نمی‌گوید می‌توانیم
دنباله‌ها را به هم الحاق کنیم، اعضایی به آن‌ها بیفزاییم، یا حتی~$d$
را از~$a_0$،~\dots،~$a_n$ با توابع تعریف‌شدنی از راه ترکیب و
کمینه‌سازی منظم \emph{محاسبه کنیم}. فقط می‌گوید تابع «رمزگشایی»ای
وجود دارد که هر دنباله‌ای «رمزگذاری» شده است.

استفاده از نماد~$\beta$ از گودل است. باز می‌گوییم که بخش دشوار اثبات
لم، تعریف~$\beta$ای مناسب با این منابع ظاهراً محدود است؛ یعنی تنها با
ترکیب و کمینه‌سازی، هرچند اجازه داریم جمع، ضرب و~$\Char{=}$ را به کار
بریم. راه‌های گوناگونی برای اثبات این لم هست، اما یکی از پاکیزه‌ترین
راه‌ها همچنان روش اصلی گودل است که از واقعیتی در نظریهٔ اعداد موسوم
به قضیهٔ سون‌زی (که به‌طور سنتی «قضیهٔ باقی‌ماندهٔ چینی» نامیده
می‌شود) استفاده می‌کند.

\begin{defn}
دو عدد طبیعی~$a$ و~$b$ \emph{نسبت به هم اول‌اند} اگر و تنها اگر
بزرگ‌ترین مقسوم‌علیه مشترکشان~$1$ باشد؛ به بیان دیگر، هیچ مقسوم‌علیه
مشترک دیگری ندارند.
\end{defn}

\begin{defn}
اعداد طبیعی~$a$ و~$b$ \emph{به پیمانهٔ~$c$ هم‌نهشت‌اند}،
$a\equiv b\mod c$، اگر و تنها اگر $c\mid(a-b)$؛ یعنی هنگام تقسیم
بر~$c$ باقیماندهٔ یکسانی دارند.
\end{defn}

قضیهٔ سون‌زی چنین است:
\begin{thm}
فرض کنید $x_0$،~\dots،~$x_n$ دو به دو نسبت به هم اول باشند. بگذارید
$y_0$،~\dots،~$y_n$ هر اعدادی باشند. آنگاه عددی چون~$z$ وجود دارد
به‌گونه‌ای که
\begin{align*}
z & \equiv y_0 \mod x_0 \\
z & \equiv y_1 \mod x_1 \\
& \vdots  \\
z & \equiv y_n \mod x_n.
\end{align*}
\end{thm}

قضیهٔ سون‌زی را چنین به کار خواهیم برد: اگر
$x_0$،~\dots،~$x_n$ به‌ترتیب از~$y_0$،~\dots،~$y_n$ بزرگ‌تر باشند،
می‌توانیم~$z$ را رمز دنبالهٔ~$\tuple{y_0,\dots,y_n}$ بگیریم. برای
بازیابی~$y_i$ فقط کافی است~$z$ را بر~$x_i$ تقسیم و باقیمانده را
برداریم. برای استفاده از این رمزگذاری باید مقادیر مناسبی برای
$x_0$،~\dots،~$x_n$ بیابیم.

چند ملاحظه در این کار به ما کمک می‌کنند. با داشتن
$y_0$،~\dots،~$y_n$، بگذارید
\begin{align*}
j &= \max(n, y_0 + 1, \dots, y_n + 1), \\
m &= \lcm(1,\dots,j),
\end{align*}
و بگذارید
\begin{align*}
x_0 & = 1 + m \\
x_1 & = 1 + 2 \cdot m \\
x_2 & = 1 + 3 \cdot m \\
& \vdots  \\
x_n & = 1 + (n+1) \cdot m.
\end{align*}
آنگاه دو امر زیر درست‌اند:
\begin{enumerate}
\item\ollabel{rel-prime} $x_0,\dots,x_n$ نسبت به هم اول‌اند.
\item\ollabel{less} برای هر~$i$، $y_i<x_i$.
\end{enumerate}
برای دیدن درستی~\olref{rel-prime}، توجه کنید اگر $p$ عددی اول باشد و
$p\mid x_i$ و~$p\mid x_k$، آنگاه
$p\mid 1+(i+1)m$ و~$p\mid 1+(k+1)m$. اما در این صورت~$p$ تفاضل
آن‌ها را تقسیم می‌کند:
\[
(1 + (i+1)m) - (1+ (k+1)m) = (i-k) m.
\]
از آنجا که~$p$، $1+(i+1)m$ را تقسیم می‌کند، نمی‌تواند~$m$ را نیز
تقسیم کند (وگرنه تقسیم نخست باقیماندهٔ~$1$ می‌داد). پس چون~$p$،
$(i-k)m$ را تقسیم می‌کند، باید~$i-k$ را تقسیم کند. اما
$\left|i-k\right|$ حداکثر~$n$ است و~$j\geq n$ را برگزیده‌ایم؛ پس
این نتیجه می‌دهد $p\mid m$، و باز به تناقض می‌رسیم. بنابراین هیچ عدد
اولی هر دو~$x_i$ و~$x_k$ را تقسیم نمی‌کند. بند~\olref{less} آسان است:
$y_i<j\leq m<x_i$.

اکنون لم تابع~$\beta$ را ثابت کنیم. به یاد آورید می‌توانیم از~$0$،
جانشین، جمع، ضرب، $\Char{=}$، افکنش‌ها و هر تابعی استفاده کنیم که
با ترکیب و کمینه‌سازیِ اعمال‌شده بر توابع منظم از آن‌ها تعریف شده
باشد. همچنین اگر تابع مشخصهٔ رابطه‌ای چنین تعریف‌شدنی باشد، می‌توانیم
آن رابطه را به کار بریم. مانند پیش می‌توانیم نشان دهیم این روابط نسبت
به ترکیب‌های بولی و کمیت‌گذاری کراندار بسته‌اند؛ برای نمونه:
\begin{align*}
\fn{not}(x) & \defis \Char{=}(x,0)\\
\bmin{x \leq z}{R(x,y)} & \defis \umin{x}{(R(x,y) \lor x = z)}\\
\bexists{x \leq z}{R(x,y)} & \defiff R(\bmin{x \leq z}{R(x,y)}, y).
\end{align*}
سپس می‌توانیم نشان دهیم همهٔ موارد زیر نیز بدون بازگشت اولیه
تعریف‌شدنی‌اند:
\begin{enumerate}
\item تابع زوج‌ساز
$J(x,y)=\frac{1}{2}[(x+y)(x+y+1)]+x$؛
% maybe explain more what is going on here, a bit confusing.
\item توابع افکنش
\begin{align*}
K(z) & = \bmin{x \leq z}{\bexists{y \leq z}{z = J(x,y)}},\\
L(z) & = \bmin{y \leq z}{\bexists{x \leq z}{z = J(x,y)}};
\end{align*}
\item رابطهٔ کوچک‌تر از~$x<y$؛
\item رابطهٔ بخش‌پذیری~$x\mid y$؛
% \item $x \tsub y$
% \item $\fn{Prime}(x)$
% \item Assuming $p$ is prime, the relation ``$x$ is a power of $p$'':
% \[
% \bforall{y \leq x}{(y \mid x \lif y = 1 \lor y = x)}.
% \]
\item تابع~$\fn{rem}(x,y)$ که باقیماندهٔ تقسیم~$y$ بر~$x$ را
  بازمی‌گرداند.
\end{enumerate}
اکنون تعریف کنید:
\begin{align*}
\beta^*(d_0,d_1,i) & = \fn{rem}(1+(i+1) d_1,d_0) \text{ و}\\
\beta(d,i) & = \beta^*(K(d),L(d),i).
\end{align*}
این همان تابعی است که می‌خواهیم. با داشتن
$a_0$،~\dots،~$a_n$ مانند بالا، بگذارید
\[
j = \max(n,a_0+1,\dots,a_n+1),
\]
و بگذارید $d_1=\lcm(1,\dots,j)$. بنا بر~\olref{rel-prime} در بالا
می‌دانیم $1+d_1$، $1+2d_1$،~\dots، $1+(n+1)d_1$ نسبت به هم اول‌اند؛
و بنا بر~\olref{less} همگی به‌ترتیب از~$a_0$،~\dots،~$a_n$ بزرگ‌ترند.
بنا بر قضیهٔ سون‌زی مقداری چون~$d_0$ وجود دارد به‌گونه‌ای که برای
هر~$i$،
\[
d_0 \equiv a_i \mod (1+(i+1)d_1)
\]
و بنابراین (چون~$d_1$ از~$a_i$ بزرگ‌تر است)،
\[
a_i = \fn{rem}(1+(i+1)d_1,d_0).
\]
بگذارید $d=J(d_0,d_1)$. آنگاه برای هر~$i\le n$ داریم:
\begin{align*}
\beta(d,i) & =  \beta^*(d_0,d_1,i) \\
& =  \fn{rem}(1+(i+1) d_1,d_0) \\
& =  a_i.
\end{align*}
این همان چیزی است که لازم داشتیم و اثبات لم تابع~$\beta$ را کامل
می‌کند.

\begin{prob}
  نشان دهید روابط~$x<y$ و~$x\mid y$ و تابع~$\fn{rem}(x,y)$ را
  می‌توان بدون بازگشت اولیه تعریف کرد. می‌توانید از~$0$، جانشین، جمع،
  ضرب، $\Char{=}$، افکنش‌ها و کمینه‌سازی و کمیت‌گذاری کراندار استفاده
  کنید.
\end{prob}

\end{document}
